\documentclass[12pt,a4paper]{article}

\usepackage{geometry}
\geometry{margin=1in}
\usepackage{array}
\usepackage{booktabs}
\usepackage{amsmath}
\usepackage{float}

\title{INFS4203/7203 Data Mining \\ Tutorial Week 2 - Introduction to Classification \\ Semester 2, 2025}
\date{}

\begin{document}
\maketitle

\section*{Question 1/4: Imputation of missing feature values and feature value normalization}

Consider a dataset in Table 1.

1. Perform missing value imputation by all-class and class-specific values!

2. Perform max-min feature normalization!

3. Perform z-score feature normalization!

4. (Extra) Does it make sense to apply normalization on categorical (nominal) features?
Justify your opinions!

\begin{table}[H]
  \centering
  \caption{A dataset of 6 data points. The right most column is the target classification label.}
  \begin{tabular}{c c c c c c}
    \toprule
    Index & Age  & Home Distance & Browser & Caffeine Level & Hair Style \\
    \midrule
    1     & 17   & 100           & Chrome  & High           & Long       \\
    2     & 20   & 50            & Chrome  & Medium         & Long       \\
    3     & null & 750           & Firefox & High           & Long       \\
    4     & 27   & 1000          & Chrome  & Low            & Short      \\
    5     & 26   & 9500          & Safari  & Low            & Short      \\
    6     & 24   & 7000          & Firefox & null           & Short      \\
    \bottomrule
  \end{tabular}
\end{table}

\textbf{Answer:}

\paragraph{For Age}
\begin{description}
  \item[Numerical:] Fill with the average value of the feature.
\end{description}

\begin{itemize}
  \item All-class: $(17 + 20 + 27 + 26 + 24) / 5 = 22.8$
  \item Class-specific: $(17 + 20) / 2 = 18.5$
\end{itemize}

\paragraph{For Caffeine Level}
\begin{description}
  \item[Nominal:] Fill with the most common feature value
\end{description}
\begin{itemize}
  \item All-class: either High or Low
  \item Class-specific: Low
\end{itemize}

\paragraph{Max-min normalization}
\begin{equation}
  x' = \frac{x - x_{\min}}{x_{\max} - x_{\min}}
\end{equation}

\begin{itemize}
  \item Age: $\max = 27, \min = 17$
        \begin{align*}
           & (17-17)/(27-17) = 0      \\
           & (20-17)/(27-17) = 0.3    \\
           & (18.5-17)/(27-17) = 0.15 \\
           & (27-17)/(27-17) = 1      \\
           & (26-17)/(27-17) = 0.9    \\
           & (24-17)/(27-17) = 0.7
        \end{align*}
  \item Home Distance: $\max = 9500, \min = 50$
        \begin{align*}
           & (100-50)/(9500-50) = 0.0053  \\
           & (50-50)/(9500-50) = 0        \\
           & (750-50)/(9500-50) = 0.0741  \\
           & (1000-50)/(9500-50) = 0.1005 \\
           & (9500-50)/(9500-50) = 1      \\
           & (7000-50)/(9500-50) = 0.7354
        \end{align*}
\end{itemize}

\paragraph{Z-Score normalization:}
\begin{equation}
  x' = \frac{x - \mu}{\sigma}
\end{equation}

\begin{itemize}
  \item Mean:
        \begin{equation}
          \mu = \frac{1}{n} \sum_{i=1}^{n} x_i
        \end{equation}
        \begin{equation*}
          \mu = \frac{17+20+18.5+27+26+24}{6} = \frac{132.5}{6} \approx 22.08
        \end{equation*}
  \item Standard deviation:
        \begin{equation}
          \sigma = \sqrt{\frac{1}{n-1} \sum_{i=1}^{n} (x_i - \mu)^2}
        \end{equation}
        \begin{equation*}
          \sigma = \sqrt{
            \frac{
              \begin{split}
                (17-22.08)^2 + (20-22.08)^2 + (18.5-22.08)^2 \\
                + (27-22.08)^2 + (26-22.08)^2 + (24-22.08)^2
              \end{split}
            }{6-1}
          }
          = \sqrt{17.24168} \approx 4.15
        \end{equation*}
  \item Calculate each feature value:
        \begin{align*}
           & (17-22.08)/4.15 = -1.22   \\
           & (20-22.08)/4.15 = -0.50   \\
           & (18.5-22.08)/4.15 = -0.86 \\
           & (27-22.08)/4.15 = 1.19    \\
           & (26-22.08)/4.15 = 0.95    \\
           & (24-22.08)/4.15 = 0.46
        \end{align*}
\end{itemize}

It does not make sense, because there are non-numeric variables representing categories without any inherent order.
\begin{itemize}
  \item Mathematically meaningless
  \item Implies ordinal relationships where none exist
\end{itemize}

% ----------------------
\section*{Question 2/4: Classifier evaluation: accuracy and F1-score on balanced datasets}

Consider the following true (ground-truth) class labels y and predicted class labels $\hat{y}_a$ and $\hat{y}_b$ from
classifiers A and B, respectively, for 6 data points (instances).

\[
  y = [1, 1, 1, 0, 0, 0],
\]
\[
  \hat{y}_a = [0, 1, 0, 0, 0, 1],
\]
\[
  \hat{y}_b = [0, 0, 1, 1, 1, 0].
\]

1. Calculate the accuracy and F1-score of Classifier A!

2. Calculate the accuracy and F1-score of Classifier B!

3. Which classifier is the best based on their accuracy scores?

4. Which classifier is the best based on their F1-scores?

\textbf{Answer:}

\subsection*{For Classifier A}
\begin{table}[H]
  \centering
  \begin{tabular}{ccc}
    \toprule
      & 1     & 0     \\
    \midrule
    1 & TP: 1 & FN: 2 \\
    0 & FP: 1 & TN: 2 \\
    \bottomrule
  \end{tabular}
\end{table}

\begin{align*}
  \text{Accuracy}  & = \frac{TP+TN}{n} = \frac{1+2}{6} = \frac{1}{2}                \\
  \text{Precision} & = \frac{TP}{TP + FP} = \frac{1}{1+1} = \frac{1}{2}             \\
  \text{Recall}    & = \frac{TP}{TP+FN} = \frac{1}{1+2} = \frac{1}{3}               \\
  \text{F1}        & = \frac{2}{\frac{1}{\text{Precision}}+\frac{1}{\text{Recall}}}
  = \frac{\frac{1}{3}}{\frac{5}{6}} = \frac{2}{5}
\end{align*}

\subsection*{For Classifier B}
\begin{table}[H]
  \centering
  \begin{tabular}{ccc}
    \toprule
      & 1     & 0     \\
    \midrule
    1 & TP: 1 & FN: 2 \\
    0 & FP: 2 & TN: 1 \\
    \bottomrule
  \end{tabular}
\end{table}

\begin{align*}
  \text{Accuracy}  & = \frac{TP+TN}{n} = \frac{1+1}{6} = \frac{1}{3}    \\
  \text{Precision} & = \frac{TP}{TP + FP} = \frac{1}{1+2} = \frac{1}{3} \\
  \text{Recall}    & = \frac{TP}{TP+FN} = \frac{1}{1+2} = \frac{1}{3}   \\
  \text{F1}        & = \frac{\frac{2}{9}}{\frac{2}{3}} = \frac{1}{3}
\end{align*}

\subsection*{}
A, because $\frac{1}{2} > \frac{1}{3}$

\subsection*{}
A, $\frac{2}{5} > \frac{1}{3}$

% ----------------------
\section*{Question 3/4: Classifier evaluation: accuracy and F1-score on imbalanced datasets}

1. Consider a cancer dataset, where 85 people have cancer and the remaining 201 do not. Classifier $\alpha$ always predicts that someone has cancer. Classifier $\beta$ always predicts that someone does not have cancer. Compare these two degenerated classifiers $\alpha$ and $\beta$ in terms of their accuracy and F1-scores! Which classifier is better (more preferable) for this domain?

2. Consider a city of 1 million people, where there are 10 COVID-positive individuals. A classifier predicts that there are 1000 COVID-positive individuals (including 10 individuals who are indeed COVID-positive), and the remaining 999,000 are COVID-negative. Analyze this classifier in terms of its accuracy and F1-score!

\textbf{Answer:}

\subsection*{Classifier $\alpha$}
\begin{table}[H]
  \centering
  \begin{tabular}{ccc}
    \toprule
              & cancer  & no cancer \\
    \midrule
    cancer    & TP: 85  & FN: 0     \\
    no cancer & FP: 201 & TN: 0     \\
    \bottomrule
  \end{tabular}
\end{table}

\begin{align*}
  \text{accuracy}  & = \frac{85}{286} \approx 0.29     \\
  \text{precision} & = \frac{85}{85+201} \approx 0.297 \\
  \text{recall}    & = 1                               \\
  \text{F1-score}  & = \frac{2}{0.297^{-1}+1^{-1}}
\end{align*}

\subsection*{Classifier $\beta$}
\begin{table}[H]
  \centering
  \begin{tabular}{ccc}
    \toprule
              & cancer & no cancer \\
    \midrule
    cancer    & TP: 0  & FN: 85    \\
    no cancer & FP: 0  & TN: 201   \\
    \bottomrule
  \end{tabular}
\end{table}

\begin{align*}
  \text{accuracy}  & = \frac{201}{286} \approx 0.703 \\
  \text{precision} & = \text{undefined}              \\
  \text{recall}    & = 0                             \\
  \text{F1-score}  & = \text{undefined}
\end{align*}

\subsection*{COVID example}
\begin{table}[H]
  \centering
  \begin{tabular}{ccc}
    \toprule
             & Positive & Negative    \\
    \midrule
    Positive & TP: 10   & FN: 0       \\
    Negative & FP: 990  & TN: 999,000 \\
    \bottomrule
  \end{tabular}
\end{table}

\section*{(Extra) Question 4/4: Classification pipeline using cross-validation (CV)}

Re-arrange the following 20 steps into a correct order of a classification pipeline with 3-fold CV.
There may be multiple correct orders out of 20! possible orders.

The nomenclature is as follows.

\begin{itemize}
  \item c($\cdot$): a trained classifier model with some arguments inside the brackets.
  \item a($\cdot$): an accuracy score with some arguments inside the brackets.
  \item mm: max-min normalization technique or the corresponding result
  \item zs: z-score normalization technique or the corresponding result
  \item avg: average, the averaging operation or the corresponding result
\end{itemize}

\begin{enumerate}
  \item Validate the trained classifier c(mm, 1, 3) using fold-2 to obtain its accuracy a(mm, 2).
  \item Train a classifier with max-min normalization using fold-1 and fold-2. Call the resulting
        trained classifier as c(mm, 1, 2).
  \item Validate the trained classifier c(mm, 1, 2) using fold-3 to obtain its accuracy a(mm, 3).
  \item Validate the trained classifier c(zs, 1, 3) using fold-2 to obtain its accuracy a(zs, 2).
  \item Train a classifier with the best hyperparameter value on training set (900 data points). Call
        the resulting trained classifier as c(zs). The training accuracy a(zs, train) may be obtained.
  \item Train a classifier with max-min normalization using fold-2 and fold-3. Call the resulting
        trained classifier as c(mm, 2, 3).
  \item Compare the average validation accuracy: a(zs, avg) versus a(mm, avg).
        Say, a(zs, avg) $>$ a(mm, avg). Thus, the best hyperparameter value is z-score normalization.
  \item Validate the trained classifier c(mm, 2, 3) using fold-1 to obtain its accuracy a(mm, 1).
  \item Split the dataset into training and testing sets. There 1000 data points in total, therefore 900
        training data points and 100 testing data points.
  \item Train a classifier with max-min normalization using fold-1 and fold-3. Call the resulting
        trained classifier as c(mm, 1, 3).
  \item Average the validation accuracy of 3 classifiers with max-min normalization.
        That is, a(mm, avg) = [a(mm, 1) + a(mm, 2) + a(mm, 3)]/3.
  \item Test the trained classifier c(zs) on testing set (100 data points) to obtain its the test accuracy
        a(zs, test). This a(zs, test) indicates the (approximate) accuracy on unseen data.
  \item Validate the trained classifier c(zs, 1, 2) using fold-3 to obtain its accuracy a(zs, 3).
  \item Train a classifier with z-score normalization using fold-2 and fold-3. Call the resulting trained
        classifier as c(zs, 2, 3).
  \item Split the training set (900 data points) into 3 folds equally. Those 3 folds are called fold-1,
        fold-2 and fold-3. There are 300 data points per fold.
  \item Validate the trained classifier c(zs, 2, 3) using fold-1 to obtain its accuracy a(zs, 1).
  \item Determine a set of hyperparameter values to tune. Here, there is one hyperparameter, i.e.
        the normalization types, namely max-min versus z-score normalizations.
  \item Train a classifier with z-score normalization using fold-1 and fold-3. Call the resulting trained
        classifier as c(zs, 1, 3).
  \item Train a classifier with z-score normalization using fold-1 and fold-2. Call the resulting trained
        classifier as c(zs, 1, 2).
  \item Average the validation accuracy of 3 classifiers with z-score normalization.
        That is, a(zs, avg) = [a(zs, 1) + a(zs, 2) + a(zs, 3)]/3.
\end{enumerate}


\end{document}
